% relatedwork.tex -- Related Work

Our work lies at the intersection of associative memory, energy-based modeling, and Langevin sampling. We review each strand and identify where existing approaches stop short of the connection we propose.

\paragraph{Hopfield networks and associative memory.}
Hopfield~\cite{hopfieldNeuralNetworksPhysical1982} introduced a recurrent network of $d$ binary neurons whose asynchronous dynamics minimize the quadratic energy $E(\mathbf{s})=-\frac{1}{2}\mathbf{s}^{\top}\mathbf{W}\mathbf{s}$, with $\mathbf{W}=\frac{1}{d}\sum_{i}\mathbf{m}_i\mathbf{m}_i^{\top}$ the Hebbian weight matrix. Stored patterns correspond to local minima of $E$, so retrieval proceeds by gradient descent from an initial probe. When the $K$ patterns are uncorrelated random binary vectors, Amit et al.~\cite{amitStoringInfiniteNumbers1985} showed via a replica-symmetric mean-field analysis that reliable retrieval is possible only when the load ratio $K/d < \alpha_c \approx 0.138$ (where $\alpha_c$ denotes the critical load ratio, not to be confused with the step size $\alpha$ used later); beyond this threshold the energy landscape becomes dominated by spurious mixtures. Folli et al.~\cite{folliMaximumStorageCapacity2017} later gave exact finite-$d$ expressions for the bit-error probability, confirming the $\alpha_c$ transition and showing that the number of retrieval errors grows with $K/d$ as cross-talk noise overwhelms the coherent signal.

The modern Hopfield network program, initiated by Krotov and Hopfield~\cite{krotovDenseAssociativeMemory2016}, addresses these shortcomings by replacing the quadratic interaction with higher-order (polynomial or exponential) interaction functions. Dense associative memories with degree-$n$ polynomial energy store on the order of $d^{n-1}$ patterns, breaking the linear capacity barrier by trading pairwise for higher-order synaptic interactions, and the resulting retrieval dynamics map onto deep networks with rectified polynomial activations. At the top of this hierarchy sits the log-sum-exp (lse) energy, which achieves storage capacity that grows \emph{exponentially} in the pattern dimension $d$~\cite{ramsauerHopfieldNetworksAll2021}. Ramsauer et al.~\cite{ramsauerHopfieldNetworksAll2021} proved that the continuous-state retrieval map associated with the lse energy is equivalent to transformer-style attention: each attention head performs one step of gradient descent on the energy, and the fixed points of the dynamics (global, metastable, or single-pattern) correspond to different retrieval regimes. All of these works, however, focus exclusively on \emph{retrieval}, that is, the convergence of iterates to stored patterns or their convex combinations. Sampling from the Boltzmann distribution defined by the same energy, which would yield a generative mechanism, is not explored. Our work fills this gap by applying Langevin dynamics to the modern Hopfield energy, converting the deterministic retrieval map into a stochastic sampler. The idea that a single energy function supports both retrieval and generation also dates to the Boltzmann machine~\cite{ackeleyLearningAlgorithmBoltzmann1985}, which shares the quadratic energy of the Hopfield network but samples from it via Gibbs dynamics instead of minimizing it; Sejnowski~\cite{sejnowskiHigherOrderBoltzmann1986} extended this to higher-order interactions, corresponding to the polynomial energies later studied by Krotov and Hopfield~\cite{krotovDenseAssociativeMemory2016}. However, these classical models operate on binary states with discrete Gibbs sampling, and our formulation inherits the same duality in the continuous, modern Hopfield setting where Langevin dynamics scales naturally to high dimensions.

\paragraph{Energy-based, score-based and diffusion models.}
Energy-based models (EBMs) define a probability distribution $p(\mathbf{x})\propto\exp(-E(\mathbf{x}))$ for an energy function $E$ parameterized by a neural network~\cite{lecunTutorialEnergyBased2006}. Training typically proceeds by contrastive divergence~\cite{hintonTrainingProductsExperts2002}, score matching~\cite{hyvarinen2005estimation}, or noise contrastive estimation~\cite{gutmannHyvaerinenNoiseContrastive2010,songHowTrainEnergyBased2021}. The energy in a generic EBM is a black-box neural network, so its gradient (the score function) must be estimated or approximated. By contrast, the modern Hopfield energy has a \emph{closed-form} gradient (identity minus softmax attention) requiring no auxiliary score network. Moreover, $\mathbf{X}$ is given as data rather than learned, so no training loop is needed, and the log-sum-exp structure provides analytic guarantees (smoothness, Lipschitz gradients, quadratic confinement) that generic EBMs achieve only with careful architectural design. The Energy Transformer~\cite{hooverEnergyTransformer2024} shows that Hopfield-type energies improve discriminative tasks (graph anomaly detection, image completion) but does not sample from the Boltzmann distribution; our work uses the same energy class for generation.

Song and Ermon~\cite{songGenerativeModelingEstimating2019} introduced score-based generative modeling, in which a neural network is trained to approximate $\nabla\log p(\mathbf{x})$ at multiple noise levels, and samples are drawn by annealed Langevin dynamics. Song et al.~\cite{songScoreBasedGenerativeModeling2021} unified this framework with diffusion probabilistic models~\cite{hoDenoisingDiffusionProbabilistic2020,sohlDicksteinDeepUnsupervisedLearning2015} by showing that both arise from forward and reverse-time stochastic differential equations. These methods achieve remarkable sample quality but rely on a learned score network trained on data; our score function is \emph{exact} and \emph{analytic}, $\nabla\log p_\beta(\boldsymbol{\xi}) = \beta(\mathbf{T}(\boldsymbol{\xi})-\boldsymbol{\xi})$, requiring no training, at the cost of being restricted to the Boltzmann distribution of a fixed memory $\mathbf{X}$, a restriction that permits convergence guarantees unavailable in the learned-score setting. Normalizing flows~\cite{rezendeMohamed2015} construct flexible generative models via learned invertible transformations; like diffusion models, they require training data and a learned bijection, and do not expose a closed-form score function over a fixed memory. Several works have introduced stochasticity into the attention mechanism itself: Deng et al.~\cite{dengLatentAlignment2018} replace the softmax with a latent categorical or Gaussian variable, trained via variational inference, to model alignment uncertainty in sequence-to-sequence tasks. These approaches treat attention weights as random variables to be inferred, whereas our stochasticity arises from Langevin dynamics on a fixed Hopfield energy: the temperature parameter $1/\beta$ controls the noise level, no variational objective is optimized, and the goal is memory-based generation rather than alignment uncertainty.